\section{Experimental Design and Setting Selection}\label{sec:attack:experiments}

\paragraph{Corpora and pools.} The corpora are Corpus A and legal128w of Chapter~\ref{ch:hub}, each with its own frozen hub.
All training data come from the \vf{40{,}000} training prompts of a corpus. % src: results/v7_2/main/manifest.json (training_pools); results/v7_2/main/anchors-linear.json (pool_size 40000)
The training prompts are divided once, at random, into two halves of \vf{20{,}000}: one half supplies the unpaired pool of the target model, as hidden states, and the other half supplies the unpaired pool of the hub, as hub outputs of OLMo.
A leak draws its paired sequences at random from all \vf{40{,}000} training prompts. % src: src/v7_2/anchors_linear.py

\paragraph{The 90-row leak.} The main experiment gives the attacker \vf{90} rows: the first nine rows of each of ten prompts. % src: w46/plan.md (leak span9); src/v7_2/anchors_linear.py (span_start_pairs); mixed-w46-std.json anchor_counts (prompts 10, leak_rows 90)
The leak has the same form on both corpora, although a Corpus A prompt averages nine words and a legal128w prompt \vf{61} (Sections~\ref{sec:hub:eval} and~\ref{sec:hub:legal}).
Ten Corpus A prompts are thus leaked almost whole, and ten legal128w prompts in their opening words.

\paragraph{Decoding and score.} Evaluation follows Section~\ref{sec:hub:eval}: each of the three decoders generates greedily from the projected output of the map, without the ground-truth prefix, up to the memory length plus \vf{32} tokens (plus \vf{8} in the Corpus A experiments other than the main comparison). % src: configs/v7_cap32.yaml; configs/v7_legal128w.yaml (cap_margin 32); configs/v7.yaml (cap_margin 8)
The score is word-F1, averaged over prompts and over the three decoders. % src: src/v7_2/f1.py; src/v7_2/factorial.py
In this chapter, word-F1 without a decoder's name denotes that three-decoder mean.

\paragraph{Control and hub reference.} The mismatched-prompt control gives each decoder the mapped states of another prompt, the one whose number of rows is nearest, and scores the output against the original prompt. % src: src/v6_4/controls.py (length_matched)
The hub reference is the word-F1 that the three decoders reach from the hub outputs of OLMo for the selection prompts, without any map: \vf{0.946} on Corpus A and \vf{0.982} on legal128w. % src: results/v7_2/main/factorial.json controls.true_olmo_code.macro_f1 (0.9459); results/legal128w/v7_2/main/mixed-v7_9-s0.json setting.corpus.g2.macro_f1 (0.9824)
It is a reference, not an upper bound: the target model and OLMo tokenize a prompt differently.

%% BIG CHANGE: polish; reword or relocate this passage for precision and flow.
\paragraph{Prompt sets.} Evaluation separates validation prompts into two disjoint sets, neither of which overlaps the training prompts.
The selection prompts are \vf{500} prompts drawn at random from the \vf{5{,}000} validation prompts of a corpus; they selected every setting of the attack. % src: v7_run.py (subset_seed 20260903, dev_prompts 500)
%% BIG CHANGE: separate the next explanatory step.

The held-out prompts are the other \vf{4{,}500} validation prompts; they selected no setting of the attack, although all \vf{5{,}000} validation prompts selected the hub checkpoint in Chapter~\ref{ch:hub}. % src: src/v7_2/context.py; results/v7_2/main/eval/confirmation/prompt_ids.json
%% BIG CHANGE: polish; reword or relocate this passage for precision and flow.
The test split of each corpus has not been evaluated.

\paragraph{Teacher-forced scores.} Two further scores are computed under teacher forcing, where the decoder receives the true preceding tokens.
The excess cross-entropy of a fit is the decoder's cross-entropy on the true next tokens under the mapped states, minus its cross-entropy under the hub outputs of OLMo.
The KL divergence measures the difference between the decoder's next-token distributions under the hub outputs of OLMo and under the mapped states.
Both are summed over the tokens of a word and reported in nats per content word, pooled over the three decoders. % src: teacher_forced_run.py; teacher_forced_summary.py; results/v7_2/main/framework/teacher_forced/summary.json (definition)

\paragraph{Experiments and their settings.} Table~\ref{tab:attack:settings} lists the experiments.
%% BIG CHANGE: polish; reorganize explanation without changing the reported evidence.
Each experiment addresses a separate question under the settings in its row.
Differences between rows in supervision, fitting, and evaluation prevent their scores from forming a single numerical progression.

\begin{table}[H]
%% BIG CHANGE: polish; reword or relocate this passage for precision and flow.
\caption{The main comparison evaluates \vf{4{,}500} held-out prompts with three seeds; each result table states the seeds of its experiment.
Settings of the hybrid fit by experiment: supervision fraction $\alpha$, temperature $\tau$ as a multiple of the scale $d$, and batch size $B$ of the matching term.
A budget counts whole paired sequences.}
\label{tab:attack:settings}
\centering
\footnotesize\setlength{\tabcolsep}{3pt}
\begin{tabular}{p{0.27\textwidth}cccp{0.22\textwidth}p{0.20\textwidth}}
\toprule
Experiment & $\alpha$ & $\tau$ & $B$ & Leak & Prompts \\
\midrule
Main comparison (H2) & \vf{0.9} & $2d$ & \vf{8{,}192} & \vf{90} rows from ten prompts & \vf{4{,}500} held-out \\ % src: configs/w46_{nq,legal128w}_std.yaml; results/*/v7_2/main/w46-confirmation-s{1,2,3}.json
Second round (exploratory) & \vf{0.95} & $2d$ & \vf{8{,}192} & budget \vf{10} or \vf{100} & \vf{450} selection \\ % src: results/*/v7_2/main/framework/attacker/*-haiku.json
Teacher-forced scores, legal128w budget curve & \vf{0.95} & $2d$ & \vf{8{,}192} & budgets \vf{1} to \vf{100} & \vf{500} selection \\ % src: results/*/v7_2/main/framework/teacher_forced/summary.json (Corpus A three seeds, legal128w seed 1)
Budget curve, Corpus A & \vf{0.95} & $2d$ & \vf{2{,}048} & budgets \vf{10} to \vf{1{,}000} & \vf{500} selection \\ % src: results/v7_2/main/mixed-v7_9-s5.json
Leak composition, second target model, decoder training & \vf{0.99} & $d$ & \vf{8{,}192} & \vf{90} rows, or budgets \vf{10} to \vf{100} & \vf{500} selection, \vf{4{,}500} held-out where stated \\ % src: results/legal128w/v7_2/main/mixed-oct_w5-s1-L90.json setting; mixed-v8_7-s1.json, mixed-v8_7-s2.json selection
\bottomrule
\end{tabular}
\end{table}

%% BIG CHANGE: polish; reword or relocate this passage for precision and flow.
The experiments retain the settings selected at their respective stages of development.
The setting of the last row was selected on legal128w alone, with whole paired sequences and one seed; the main comparison then fixed one setting for both corpora, as follows. % src: results/legal128w/v7_2/main/mixed-v8_7-s1.json selection (alpha 0.99), mixed-v8_7-s2.json selection (tau 0.06485)

%% BIG CHANGE: polish; reword or relocate this passage for precision and flow.
\paragraph{Grid search.} The main comparison fixes its setting through a grid search over the supervision fraction and the temperature, on the \vf{90}-row leak, with seed \vf{1} and the \vf{500} selection prompts of each corpus (Figure~\ref{fig:attack:grid}). % src: results/v7_2/main/mixed-w46-grid.json; results/legal128w/v7_2/main/mixed-w46-grid.json; w46-selection.json
The rule, written before the grid was run, selects the cell with the highest mean word-F1 over the two corpora. % src: w46/plan.md; w46-selection.json rule
The selected cell is $\alpha = 0.9$, $\tau = 2d$, with a two-corpus mean of \vf{0.534}. % src: w46-selection.json best (0.5336)
It is the best cell on Corpus A (\vf{0.768}) and not on legal128w, where it scores \vf{0.299} against \vf{0.314} for the best cell ($\alpha = 0.8$, $\tau = 2d$). % src: mixed-w46-grid.json seeds.1.cells (0.7679; legal128w 0.2993, 0.3144)
On both corpora word-F1 falls on either side of $\tau = 2d$ for $\alpha \le 0.95$, and several cells with a low temperature score below the supervised fit.
%% BIG CHANGE: polish; reword or relocate this passage for precision and flow.
The matching term therefore helps only within a range of temperatures, and selecting that range requires prompts with known text.

\begin{figure}[ht]
\centering
\includegraphics[width=0.95\textwidth]{ch4/grid_search}
\caption{One setting, $\alpha = 0.9$ and $\tau = 2d$ (outlined), has the highest mean word-F1 over the two corpora.
Word-F1 of the hybrid fit by supervision fraction (rows) and temperature (columns), on the \vf{90}-row leak, seed \vf{1}, \vf{500} selection prompts per corpus, $B = 8{,}192$.
The titles give the supervised fit on the same leak.}
\label{fig:attack:grid}
\end{figure}
